\documentclass{beamer}

\usetheme{Madrid}
\usecolortheme{dolphin}

\title{Prime Time}
\author{Siddharth \& Amit}
\date{\today}

\newcommand{\Z}{\mathbb{Z}}
\newcommand{\bld}[1]{\textbf{#1}}
\newcommand{\itl}[1]{\textit{#1}}
\newcommand{\N}{\mathbb{N}}
\newcommand{\Q}{\mathbb{Q}}
\newcommand{\R}{\mathbb{R}}
\newcommand{\eps}{\varepsilon}
\newcommand{\poly}{\mathrm{poly}}
\newcommand{\polylog}{\mathrm{polylog}}
\newcommand{\bigO}{\mathcal{O}}

\begin{document}

\frame{\titlepage}

\section{The Problem of PRIMES}

\begin{frame}{Mathematical Definition}
    Let $n \in \mathbb{Z}$ with $n > 1$. Determine whether $n$ is \textbf{prime}.
\end{frame}

\begin{frame}{PRIMES as a Decision Problem}

\textbf{Input:} An integer $n > 1$

\textbf{Output:}
\begin{itemize}
    \item YES if $n$ is prime
    \item NO otherwise
\end{itemize}

\end{frame}

\begin{frame}{Language Formulation}

We define the language:

\[
\textbf{PRIMES} = \{\, n \in \mathbb{Z} \mid n > 1 \text{ and } n \text{ is prime} \,\}
\]

Equivalently (binary encoding):

\[
\textbf{PRIMES} = \{\, \langle n \rangle \mid n \text{ is prime} \,\}
\]

where $\langle n \rangle$ is the binary representation of $n$.

\end{frame}

\begin{frame}{Computational Question}

\begin{block}{Goal}
Design an algorithm that decides whether:
\[
n \in \textbf{PRIMES}
\]
\end{block}

\pause

\begin{itemize}
    \item How efficiently can we decide this?
    \item Is $\textbf{PRIMES} \in \mathbf{P}$?
\end{itemize}

\end{frame}

%------------------------------------------------

\begin{frame}{Trial Division is a Lie}

\textbf{Idea:} Check if $n$ is divisible by any integer up to $\sqrt{n}$

\[
\text{Time} = O(\sqrt{n})
\]

\pause

Let $k = \log n$ (input size)

\[
\sqrt{n} = n^{1/2} = (2^k)^{1/2} = 2^{k/2}
\]

\pause

\begin{block}{Key Insight}
$O(\sqrt{n}) = O(2^{k/2})$ is \textbf{exponential in input size}
\end{block}

\end{frame}

%------------------------------------------------

\begin{frame}{Why This is Bad}

\begin{itemize}
    \item Input size is $k = \log n$
    \item Algorithm runs in $O(2^{k/2})$
\end{itemize}

\pause

\begin{block}{Conclusion}
Trial division is \textbf{EXP time}, not polynomial
\end{block}

\pause

\begin{center}
\Large
\textit{Can we do better?}
\end{center}

\end{frame}

%------------------------------------------------

\begin{frame}{The Core Question}

\begin{block}{Computational Problem}
Does there exist a deterministic algorithm that decides primality in time polynomial in $k$?
\end{block}

\pause

\begin{center}
\Large
\[
\textbf{PRIMES} \stackrel{?}{\in} \mathbf{P}
\]
\end{center}

\end{frame}

%------------------------------------------------

\section{Improving Trial Division}

\begin{frame}{Wheel Factorization (Idea)}

Instead of testing all numbers:

\begin{itemize}
    \item First check divisibility by small primes
    \item Skip numbers that are obviously composite
\end{itemize}

\pause

Example:
\begin{itemize}
    \item Check $2$
    \item Then only test odd numbers
\end{itemize}

\end{frame}

%------------------------------------------------

\begin{frame}{Primorials}

\[
c\# = \prod_{p \le c} p
\]

\pause

\begin{itemize}
    \item Product of all primes $\le c$
    \item Used to eliminate large chunks of candidates
\end{itemize}

\end{frame}

%------------------------------------------------

\begin{frame}{Key Observation}

All integers can be written as:

\[
x = c\# \cdot k + i
\]

\pause

Only values of $i$ that are \textbf{coprime to $c\#$} can yield primes

\end{frame}

%------------------------------------------------

\begin{frame}{Example: $6\# = 30$}

\[
30 = 2 \cdot 3 \cdot 5
\]

All $x \ge 30$ can be written as:
\[
x = 30k + i
\]

Eliminate residues divisible by $2,3,5$

\[
i \in \{1,7,11,13,17,19,23,29\}
\]

\end{frame}

%------------------------------------------------

\begin{frame}{Search Space Reduction}

\begin{itemize}
    \item Total residues: $30$
    \item Valid candidates: $8$
\end{itemize}

\[
\text{Reduction} = 1 - \frac{8}{30} = 76.67\%
\]

\begin{block}{Important}
These are \textbf{candidates}, not guaranteed primes
\end{block}

\end{frame}

%------------------------------------------------

\begin{frame}{Limitations}

\begin{itemize}
    \item Still need to test divisibility
    \item Must check primes $\le \sqrt{n}$
\end{itemize}

\begin{block}{Conclusion}
Wheel factorization reduces constants, \\
but does \textbf{not} change exponential nature
\end{block}

\end{frame}

%------------------------------------------------

\begin{frame}{Towards Better Methods}

\begin{itemize}
    \item Repeated elimination ideas lead to:
    \item \textbf{Sieve of Eratosthenes}
\end{itemize}

\begin{block}{Big Picture}
We need fundamentally different ideas \\
to get polynomial time
\end{block}

\end{frame}

\end{document}